\uuid{czPq}
\titre{Étude de suite}
\auteur{bodin ; hanani ; mzari}
\organisation{exo7}
\licence{CC BY-NC-SA 4.0}
\datecreate{2020-01-20}
\niveau{L1}
\module{Analyse}
\chapitre{Suite}
\sousChapitre{Autre}
\difficulte{3}
\theme{Suites, étude de suite}
\exo7qcmid{lille-1.suites.027}

\contenu{
\begin{question}
\qclassification{121.00}
\qtags{motcle=étude de suite}
Soit $(u_n)_{n\geq 0}$ la suite définie par $u_0\geq 0$ et $\displaystyle u_{n+1}= \ln(1+u_n)$. Que peut-on en déduire ?
\begin{answers}
    \bad{Une telle suite $(u_n)$ n'existe pas.}
    \good{$\forall n\in \Nn^*$, $u_n \geq 0$, et $(u_n)$ est décroissante}
    \bad{$\displaystyle \lim _{n\to +\infty}u_n=+\infty$}
    \good{$\displaystyle \lim _{n\to +\infty}u_n=0$}
\end{answers}
\begin{explanations}
On vérifie par récurrence que $\displaystyle 0\leq u_n$ pour tout $n\geq 0$. Donc la suite $(u_n)$ est bien définie. On vérifie aussi que $\ln (1+x)\leq x$ pour tout réel $x\geq 0$. En particulier,
$$u_{n+1}=\ln (1+u_n)\leq u_n.$$
Donc $(u_n)$ est décroissante. Étant à la fois décroissante est minorée, la suite $(u_n)$ est convergente et sa limite est l'unique solution de l'équation $x=\ln (1+x)$. Soit $\displaystyle \lim _{n\to +\infty}u_n=0$.
\end{explanations}
\end{question}
}
